\documentclass[a4paper,11pt]{article}
\usepackage{exam-style-842-unit}
\examinehead{842 信息化概论 · 第 1 单元练习册}{练习二 参考答案与解析}
\begin{document}

\unitanswerscover{842 信息化概论 · 单元练习册}{第 1 单元参考答案}{练习二（信息社会与发展战略）的答案与解析；每道判断题均给出改正后的正确表述。}

\answerpart{一、单项选择题（10 题，每题 3 分，共 30 分）}

\q B。\begin{analysis}信息社会也叫信息化社会、知识社会、网络社会、虚拟社会、后工业社会等，其战略资源是信息和知识（教材 1.2.1，第 5—6 页）。A 资本与机器是工业社会的关键资源；C 土地与劳动力是农业社会的关键资源；D 能源与材料属于物质性资源。\end{analysis}

\q C。\begin{analysis}信息社会起主导作用的是知识密集型产业（教材 1.2.1）。A 劳动密集型产业是农业社会及工业化初期的典型形态；B 资本密集型产业是工业社会的主导形态；D 资源密集型产业依赖自然资源投入，均不是信息社会的主导产业。\end{analysis}

\q B。\begin{analysis}这是教材 1.2.1 对信息化的定义：信息化是充分利用信息技术，开发利用信息资源，促进信息交流和知识共享，提高经济增长质量，推动经济社会发展转型的历史进程。A 数字化是将模拟信号通过抽样、量化、编码等环节变换成数字码元；C 工业化是机器大生产逐步取代手工劳动的过程；D 城镇化是人口向城镇聚集、城镇生产生活方式扩散的过程。\end{analysis}

\q A。\begin{analysis}1963 年日本社会学家梅棹忠夫在《信息产业论》中首次提出“信息社会”概念；1979 年贝尔认为“信息社会”的概念比“后工业社会”更确切；托夫勒 1980 年在《第三次浪潮》中系统阐述第三次浪潮思想；维纳是控制论创始人。（教材 1.2.1）\end{analysis}

\q D。\begin{analysis}信息社会的基本特征有 4 条：知识型经济、网络化社会、数字化生活、服务性政府（教材 1.2.1，第 6—7 页）。机械化生产是工业社会的特征，不属于信息社会的基本特征。\end{analysis}

\q A。\begin{analysis}信息社会发展阶段按信息社会指数（ISI）划分为：起步期（ISI 值 0.3 以下）、转型期（0.3—0.6）、初级阶段（0.6—0.8）、中级阶段（0.8—0.9）、高级阶段（0.9 以上）（教材 1.2.2，第 7—8 页）。\end{analysis}

\q B。\begin{analysis}电子政务指政府运用信息技术优化管理与服务流程，提高行政效率、增强公共服务能力，向社会提供规范、透明、高效的管理和服务。A 政府网上采购只是电子政务的一种具体应用；C 电子商务的主体是企业、对象是商务活动，与电子政务不同；D 电子政务不仅面向内部办公，更重要的是面向社会和公众服务。电子政务被《2006—2020 年国家信息化发展战略》列为重点任务之一（教材 1.5.2）。\end{analysis}

\q B。\begin{analysis}计算思维是概念化思维，不是程序化思维；是人的思维，不是计算机的思维；是数学和工程互补融合的思维，不是纯数学性思维；是面向所有人的思维，而不仅仅是计算机科学家的思维（教材 1.4.2，第 13—14 页）。A、C、D 均与教材表述相反。\end{analysis}

\q B。\begin{analysis}从人机共生思想转到基于三元社会模式的新信息世界观，是信息技术认识转变的第 3 点：今天的信息世界是多人、多机、多物组成的动态开放网络社会，即物理世界、信息世界、人类社会组成的三元世界（教材 1.4.3）。A 人机共生是转变之前的旧思想；C 狭义工具论是把信息科技仅看作工具的狭隘认知；D 计算思维是认识转变的第 2 点。\end{analysis}

\q A。\begin{analysis}《2006—2020 年国家信息化发展战略》是我国信息化发展的纲领性文件，提出了到 2020 年的战略目标和 9 项重点（教材 1.5，第 14—15 页）。B 是国家信息中心发布的信息社会发展定量测评报告；C 是托夫勒的著作；D 是维纳的控制论专著。\end{analysis}

\answerpart{二、多项选择题（3 题，每题 4 分，共 12 分）}

\q ABC。\begin{analysis}信息化特性可归纳为“四化”和“四性”，“四化”指智能化、电子化、全球化、非群体化（教材 1.2.1）。D 错误：信息时代人们的活动更加个性化，个人之间的信息交换日益增加，非群体化才是信息化的特征，“群体化”与教材表述相反。\end{analysis}

\q ABC。\begin{analysis}“四性”指综合性、竞争性、渗透性和开放性（教材 1.2.1，第 5—6 页）。D 错误：开放性是信息化的特性之一，“封闭性”与教材表述相反。\end{analysis}

\q ABCD。\begin{analysis}信息社会的基本特征包括：知识型经济（人力资源知识化、发展方式可持续化、产业结构软化、经济水平发达）；网络化社会（信息基础设施的完备性、社会服务的包容性、社会发展的协调性）；数字化生活（生活工具数字化、生活方式数字化、生活内容数字化）；服务性政府（科学决策、公开透明、高效治理、互动参与）。四项全对。（教材 1.2.1）\end{analysis}

\answerpart{三、判断题（6 题，每题 2 分，共 12 分）}

\q √。\begin{analysis}《2006—2020 年国家信息化发展战略》提出，到 2020 年综合信息基础设施基本普及，信息技术自主创新能力显著增强，信息产业结构全面优化，国家信息安全保障水平大幅提高，国民经济和社会信息化取得明显成效等（教材 1.5.1）。\end{analysis}

\q ×。改正：信息社会是脱离工业化社会以后信息起主要作用的社会，它以信息技术为基础、以信息产业为支柱、以信息价值的生产为中心、以信息产品为标志，战略资源是信息和知识；互联网普及只是其技术基础条件之一，并不等于建成了信息社会。\begin{analysis}把信息社会等同于互联网普及的社会，忽略了信息产业、信息价值的生产、知识密集型产业等本质规定。（教材 1.2.1）\end{analysis}

\q ×。改正：信息社会的战略资源是信息和知识，起主导作用的是知识密集型产业；资本和能源是工业社会的战略资源。\begin{analysis}教材 1.2.1 明确指出：信息社会起主导作用的是知识密集型产业，战略资源是信息和知识。\end{analysis}

\q √。\begin{analysis}人类社会的四次信息技术革命依次为：语言和文字；造纸和印刷技术；电报、电话、电视及其他通信技术；电子计算机和现代通信技术在信息工作中的应用（教材 1.3.1，第 8 页）。\end{analysis}

\q √。\begin{analysis}信息化的实现要以信息的数字化为必备条件，数字化则为信息化的实现提供有力的技术支撑（教材 1.3.3）。\end{analysis}

\q ×。改正：计算思维是概念化思维而不是程序化思维，是人的思维而不是计算机的思维，是面向所有人的普适思维；编写计算机程序只是其具体应用之一，二者不能等同。\begin{analysis}计算思维强调一切皆可计算，从物理世界模拟到人类社会模拟再到智能活动都可认为是计算的某种形式，其内涵远大于程序设计。（教材 1.4.2）\end{analysis}

\answerpart{四、填空题（6 空，每空 1 分，共 6 分）}

\q （1）非群体化。\begin{analysis}信息化“四化”指智能化、电子化、全球化、非群体化（教材 1.2.1）。\end{analysis}

\q （2）信息产业。\begin{analysis}信息社会以信息技术为基础，以信息产业为支柱，以信息价值的生产为中心，以信息产品为标志（教材 1.2.1）。\end{analysis}

\q （3）造纸和印刷技术。\begin{analysis}第二次信息技术革命是造纸和印刷技术的出现，它结束了人们单纯依靠手抄、雕刻文献的时代，使知识可以大量生产、存储和流通（教材 1.3.1）。\end{analysis}

\q （4）ISI；起步。\begin{analysis}信息社会发展指数英文缩写为 ISI，其值 0.3 以下为起步期，0.3—0.6 为转型期（教材 1.2.2）。\end{analysis}

\q （5）计算思维。\begin{analysis}信息科学技术认识转变的第 2 点是从狭义工具论转到计算思维（教材 1.4.2）。\end{analysis}

\answerpart{五、简答题（4 题，每题 6 分，共 24 分）}

\q 参考答案要点：
（1）信息化是充分利用信息技术，开发利用信息资源，促进信息交流和知识共享，提高经济增长质量，推动经济社会发展转型的历史进程（教材 1.2.1，第 5 页）。
（2）信息化是一个国家由物质生产向信息生产、由工业经济向信息经济、由工业社会向信息社会转变的动态的、渐进的过程，也是一个社会经济结构不断变换的过程。
（3）“四化”：智能化、电子化、全球化、非群体化。
（4）“四性”：综合性、竞争性、渗透性、开放性。

\q 参考答案要点：
（1）信息社会也称信息化社会，是脱离工业化社会以后信息起主要作用的社会；它以信息技术为基础，以信息产业为支柱，以信息价值的生产为中心，以信息产品为标志，战略资源是信息和知识，起主导作用的是知识密集型产业（教材 1.2.1）。
（2）基本特征有 4 条：
（1）知识型经济：人力资源知识化、发展方式可持续化、产业结构软化、经济水平发达；
（2）网络化社会：信息基础设施的完备性、社会服务的包容性、社会发展的协调性；
（3）数字化生活：生活工具数字化、生活方式数字化、生活内容数字化；
（4）服务性政府：科学决策、公开透明、高效治理、互动参与。

\q 参考答案要点：
（1）阶段划分：国家信息中心 2010 年发布《走近信息社会：中国信息社会发展报告 2010》，提出信息社会发展指数（ISI）和测评体系。信息社会发展大体分为起步期（ISI 值 0.3 以下）、转型期（0.3—0.6）、初级阶段（0.6—0.8）、中级阶段（0.8—0.9）、高级阶段（0.9 以上）（教材 1.2.2）。
（2）我国发展水平：2012 年我国信息社会指数为 0.4391，比 2010 年提高 17\%，仍处于从工业社会向信息社会过渡的加速转型期。
（3）不平衡状况：2012 年东、中、西部地区信息社会指数均值分别为 0.541、0.386、0.358，发展水平最高省份的信息社会指数是最低省份的 2.5 倍；2015 年东、中、西部信息社会指数分别为 0.5489、0.3880、0.3729，东部地区比全国平均水平高 26.15\%，比中、西部地区分别高 41.5\% 和 47.20\%。这表明我国存在明显的数字鸿沟，区域发展不平衡。

\q 参考答案要点：
（1）认识转变的 4 个方面（教材 1.4，第 12—14 页）：（一）从重视信息技术的内涵转到更加重视其外延；（二）从狭义工具论转到计算思维；（三）从人机共生思想转到基于三元社会模式的新信息世界观；（四）信息科学技术重点研究方向的改变。
（2）计算思维的要点：计算思维是一种普适的思维，是每个人的基本技能，强调一切皆可计算；它是概念化思维而不是程序化思维，是人的思维而不是计算机的思维，是数学和工程互补融合的思维而不是纯数学性思维，是面向所有人的思维而不仅仅是计算机科学家的思维。
（3）重点研究方向的改变：长期以来信息科学技术研究的主要目标是提高信息器件和系统的性能，摩尔定律指引的方向主要是提高半导体器件的集成度，从而提高主频和性能；现在 CMOS 器件主频的提高已受到功耗的限制，今后发展信息技术的主要致力方向将是降低功耗、成本和体积（占地面积），提高易用性、效率和性能（教材 1.4.4）。
（4）21 世纪信息技术发展的新取向：在继续发展工程技术规模效益的同时，更加重视信息技术的多样性、开放性和个性化，更加重视信息技术惠及大众；更加重视生态和环境影响，探索对有限自然资源和无限知识资源的分享、共享和持续利用；更加重视医学及与人类健康有关的信息科学技术；更加重视信息科学的研究探索，特别是与纳米、生命、认知等科学的交叉研究；更加重视信息技术与人文艺术的结合。
（5）说明：计算机技术的具体主要发展领域在教材第 2 章展开，属本课程后续单元内容；就第 1 章而言，重点发展方向即上述“降低功耗、成本和体积，提高易用性、效率和性能”及 21 世纪新取向。

\answerpart{六、综合题（论述，1 题，共 16 分）}

\q 参考答案要点：
（1）信息技术改变农业生产（4 分）：气象、土壤、市场行情等数据经采集、传输、处理成为可共享的信息资源，测土配方施肥把土壤养分数据转化为施肥决策，病虫害在线诊断把图像与知识库比对转化为防治建议，农产品电商把市场需求信息反馈给生产端。信息技术已成为新的生产要素，推动农业生产由经验驱动向数据驱动、精准化转变。
（2）信息技术改变政府治理（4 分）：农业信息服务平台使农业部门的管理与服务由层层上报、纸质传递转向在线协同，是电子政务在农业农村领域的具体体现；信息资源共享使农户、合作社、企业、政府能够同时获得同一份数据，提高决策透明度和公共服务效率，服务性政府特征逐步显现。
（3）信息产业与信息资源共享（4 分）：平台建设带动了信息设备制造、软件开发、信息服务和数据服务等信息产业部门的发展，信息产业具有高渗透、高催化、知识智力密集的特征；信息资源共享是信息可共享性的体现，信息不因共享而损耗，共享范围越大其价值越高，因此要推动跨部门、跨层级的数据共享与业务协同。
（4）数字鸿沟与对策（4 分）：我国东、中、西部地区信息社会指数差距明显，农村地区在信息基础设施、终端普及和农民信息素质方面相对薄弱，容易形成数字鸿沟。对策包括：完善农村综合信息基础设施，提高农民信息技术应用能力，加强信息资源的开发利用，提高信息产业竞争力，建设信息安全保障体系，并把信息化与工业化、城镇化、农业现代化同步推进。这正是我国信息化发展战略的目标与重点所在。（教材 1.2.2、1.4、1.5）

\end{document}
